\chapter{The Stock Effects}
\label{ch:stock-effects}

XNA shipped five built-in shader effects, and CNA ports all five. This chapter covers the
shared \cnaclass{Effect} base and its supporting types first, then each stock effect, and
then the specific, cross-backend bug pattern their verification kept turning up --- which is
worth understanding as a pattern, not just five separate bug lists.

\section{Effect: the base, and why its bytecode constructor always throws}

\cnaclass{Effect} inherits \cnaclass{GraphicsResource} and is not copyable. Its two
constructors tell the whole story of Chapter~\ref{ch:shader-fx-gap} in miniature:
\texttt{Effect(device)} is the ordinary path every stock effect uses, while
\texttt{Effect(device, vector<bytecs> effectCode)} --- the constructor that would load a
real compiled \texttt{.fx} bytecode blob, the way a real XNA content-pipeline output does ---
\textbf{always throws \texttt{NotImplementedException}}. CNA has no equivalent of
MojoShader's bytecode parser, and Chapter~\ref{ch:shader-fx-gap} covers why that gap is
architectural rather than an oversight. \texttt{Clone()} is pure virtual and, as a
\texttt{NOXNA} deviation from FNA's garbage-collected reference semantics, returns an owning
raw pointer rather than a managed reference --- there is no GC in C\texttt{++} to hand a
shared reference back to. \texttt{Apply()} itself is \texttt{NOXNA} --- real FNA has no
public \texttt{Effect.Apply()}, only \texttt{EffectPass.Apply()} --- and the protected virtual
\texttt{FillGpuDrawParams(GpuDrawParams\&)} is the single mechanism every stock effect below
uses to hand its own parameters down to whichever backend is active; the base implementation
is a no-op, and every stock effect overrides it.

\cnaclass{EffectParameter} carries full metadata (name, semantic, row/column counts,
\cnaclass{EffectParameterClass}, \cnaclass{EffectParameterType}) plus typed
\texttt{GetValue*} / \texttt{SetValue*} accessors for every scalar, vector, matrix, quaternion,
string, and texture type XNA effects support. \cnaclass{EffectParameterCollection},
\cnaclass{EffectPassCollection}, and \cnaclass{EffectTechniqueCollection} all store their
elements behind \texttt{unique\_ptr} rather than by value, specifically so that a reference
obtained earlier --- including a captured \texttt{CurrentTechnique} --- stays valid even if a
later \texttt{Add()} reallocates the backing vector. \cnaclass{EffectPass::Apply()} throws
\texttt{InvalidOperationException} if the pass does not belong to the effect's currently
selected technique, including when \texttt{CurrentTechnique} is null --- matching FNA's own
``Applied a pass not in the current technique!'' guard, mapped onto a defined C\texttt{++}
exception rather than the null-reference crash the equivalent null case produces in C\#.

\section{GpuDrawParams: the one struct every backend reads}

Every stock effect's \texttt{FillGpuDrawParams} override populates the same struct, defined
once in the backend contract layer (Chapter~\ref{ch:backend-architecture}) and forwarded from
\texttt{Effect::FillGpuDrawParams()} to whichever backend is currently active. It carries two
ordinary texture slots plus a dedicated environment-map cube slot, diffuse/ambient/emissive/%
specular colors, three directional lights (each with direction, diffuse, and specular),
the world matrix, alpha-test parameters, Fresnel/environment-map parameters, a 72-bone
skinning palette with a weights-per-vertex count, fog parameters, a set of mode flags
(\texttt{textureEnabled}, \texttt{vertexColorEnabled}, \texttt{lightingEnabled},
\texttt{dualTexture}, \texttt{envMapping}, \texttt{skinned}, \texttt{pbr}), instancing fields,
a pointer for a custom \texttt{ShaderEffect}-compiled program, and six PBR-specific
metallic-roughness fields. Two flags are worth flagging directly: \texttt{preferPerPixelLighting}
and \texttt{specularEnabled} are both documented in the struct itself as ``ignored by every
backend except D3D9'' --- an explicitly acknowledged, tracked divergence from XNA's own
default lighting behavior on every backend except the one built to be pixel-authentic
(Chapter~\ref{ch:direct3d-backends}).

\section{BasicEffect}

\cnaclass{BasicEffect} is the largest of the five: public \texttt{World} / \texttt{View}/%
\texttt{Projection} fields, \texttt{VertexColorEnabled}, the full material color set
(\texttt{DiffuseColor}, \texttt{EmissiveColor}, \texttt{SpecularColor}, \texttt{SpecularPower},
\texttt{Alpha}, \texttt{AmbientLightColor}), \texttt{LightingEnabled}/%
\texttt{PreferPerPixelLighting}, \texttt{TextureEnabled} / \texttt{Texture}, full
\texttt{DirectionalLight0/1/2}, and \texttt{EnableDefaultLighting()}. Its 22-property audit
found and fixed several early default-value bugs (\texttt{VertexColorEnabled},
\texttt{DirectionalLight0.Enabled}, and \texttt{DirectionalLight.Direction}'s default value
were all initially wrong), then a sequence of real, per-backend rendering bugs: the
no-texture \texttt{VertexColorEnabled} toggle was ignored by all three of EasyGL, Vulkan, and
BGFX (three separate bugs, one per backend); the textured-plus-vertex-color-plus-diffuse case
dropped \texttt{DiffuseColor} entirely on EasyGL and BGFX (Vulkan was already correct); and a
shared bug across all three backends never checked \texttt{DirectionalLight0.Enabled} at all,
so a disabled light kept lighting the surface regardless. BGFX carried one additional,
much wider bug: its vertex-layout builder never declared \texttt{Normal} / \texttt{TexCoord0}
attributes for any stride except the skinned 52-byte layout --- silently breaking lit-normal
\emph{and} UV interpolation for every other stride, invisible until then because every prior
test happened to use a UV-insensitive $1\times1$ texture. \texttt{EmissiveColor} was dropped
entirely in the no-lighting path on all three backends; ambient/emissive/specular forwarding
and real specular highlights were all fixed afterward, leaving all three backends
byte-identical for the full feature composition as of the most recent audit.

\subsection{A worked example: a lit, textured cube}

\cnaclass{BasicEffect} is deliberately the effect a first 3D CNA program reaches for, and its
own real headers confirm why it reads like plain XNA rather than a wrapped shader system:
\texttt{World} and \texttt{VertexColorEnabled} are ordinary public fields (matching real XNA's
own field-not-property choice for this specific effect), while \texttt{Texture} and
\texttt{TextureEnabled} go through the usual \texttt{getXProperty} / \texttt{setXProperty} pair:

\begin{lstlisting}
BasicEffect effect(getGraphicsDeviceProperty());
effect.World = Matrix::CreateRotationY(rotationAngle);
effect.setViewProperty(camera.View());
effect.setProjectionProperty(camera.Projection());

effect.setTextureEnabledProperty(true);
effect.setTextureProperty(&cubeTexture);

effect.EnableDefaultLighting();          // sets DirectionalLight0/1/2 to XNA's own defaults
effect.DirectionalLight0.setDiffuseColorProperty(Vector3(1.0f, 0.95f, 0.9f)); // warm key light

for (EffectPass& pass : effect.getCurrentTechniqueProperty()->getPassesProperty()) {
    pass.Apply();
    getGraphicsDeviceProperty().DrawIndexedPrimitives(/* ... */);
}
\end{lstlisting}

\texttt{EnableDefaultLighting()} is worth calling before touching any individual
\texttt{DirectionalLightN} property, not after --- reading its own implementation directly
confirms it unconditionally overwrites \texttt{DirectionalLight0}'s diffuse color, direction,
specular color, and enabled flag (and likewise for lights 1 and 2) every time it runs; setting
\texttt{DirectionalLight0}'s diffuse color first and calling \texttt{EnableDefaultLighting()}
second would silently overwrite the custom color right back to the default. This is exactly the kind of ordering trap the per-backend bug list below
does \emph{not} cover, because it is a real API-usage hazard rather than a backend bug ---
\texttt{EnableDefaultLighting()}'s own real behavior, faithfully ported, not a CNA-specific
quirk.

\section{AlphaTestEffect}

\cnaclass{AlphaTestEffect} implements \texttt{IEffectMatrices} and \texttt{IEffectFog} only
--- no lighting. \texttt{AlphaFunction} (a \cnaclass{CompareFunction}, defaulting to
\texttt{Greater}) and \texttt{ReferenceAlpha} (an unclamped 0--255 int) implement the actual
alpha test; \texttt{VertexColorEnabled}, \texttt{DiffuseColor} / \texttt{Alpha}, and
\texttt{Texture} round out the surface. Its audit found zero property-default bugs (the
first test coverage for this effect written from scratch), and a full 8-value
\cnaclass{CompareFunction} sweep passed 24 of 24 cases on all three backends. One real,
still-open gap: \texttt{VertexColorEnabled} has no effect at all on Vulkan or BGFX, because
both route every alpha-test draw through one generic pipeline/shader that only declares
position and texture-coordinate inputs, no color attribute --- described in the project's
own tracking as ``a genuinely large, six-shader-file, two-backend change,'' not attempted.
EasyGL, by contrast, handles it correctly. Fog forwarding was a real, fixed bug on EasyGL
specifically for this effect --- and, in the course of fixing it, revealed a much larger
finding: fog was a \emph{total} no-op on both Vulkan and BGFX for every 3D effect, because no
shader file on either backend mentioned fog at all. That project-wide gap was closed in one
pass, adding real fog uniforms, varyings, and blend formula to every 3D shader on both
backends at once --- the same ``one shared fix closes it everywhere'' pattern already seen
for \cnaclass{SpriteBatch}'s flip bug in Chapter~\ref{ch:spritebatch}. A separate, unrelated
BGFX-only bug --- seven texture-binding call sites with no fallback-to-white-texture
else-branch --- was found and fixed at the same time, affecting every effect that binds a
possibly-null texture, not just this one.

\subsection{A worked example: a cutout-alpha chain-link fence}

The classic real use of \cnaclass{AlphaTestEffect} is a cutout texture --- foliage, a
chain-link fence, a wire mesh --- where a pixel is either fully opaque or fully discarded, with
no blended edge at all (unlike ordinary alpha \emph{blending}, which this effect deliberately
does not do). \texttt{AlphaFunction} and \texttt{ReferenceAlpha} together express the test as a
comparison against the texture's own alpha channel:

\begin{lstlisting}
AlphaTestEffect effect(getGraphicsDeviceProperty());
effect.setWorldProperty(fenceWorld);
effect.setViewProperty(camera.View());
effect.setProjectionProperty(camera.Projection());

effect.setTextureProperty(&fenceTexture);         // alpha channel: 255 = wire, 0 = gap
effect.setAlphaFunctionProperty(CompareFunction::Greater);
effect.setReferenceAlphaProperty(128);            // keep only pixels with alpha > 128

for (EffectPass& pass : effect.getCurrentTechniqueProperty()->getPassesProperty()) {
    pass.Apply();
    getGraphicsDeviceProperty().DrawIndexedPrimitives(/* ... */);
}
\end{lstlisting}

This effect's real, still-open gap from the audit above bites exactly this kind of asset if it
also carries per-vertex color (a common way to fade fence-post tint by ambient occlusion baked
into the mesh): setting \texttt{setVertexColorEnabledProperty(true)} here has zero visible
effect on Vulkan or BGFX specifically --- not a subtle blend-formula difference, but a complete
no-op, because neither backend's alpha-test shader declares a color vertex attribute at all. A
game targeting all three hardware backends should not rely on \cnaclass{AlphaTestEffect}'s
vertex color for anything visually load-bearing until that gap closes.

\section{DualTextureEffect}

Two independent texture slots (\texttt{Texture}, \texttt{Texture2}), plus
\texttt{VertexColorEnabled}, \texttt{DiffuseColor} / \texttt{Alpha}, and fog --- no lighting.
The audit's headline finding is a real formula bug present on all three backends at once:
FNA's real shader doubles the first texture's RGB value (\texttt{color.rgb *= 2}) before
multiplying it against the overlay texture, and none of CNA's three backend shaders
implemented that factor at all --- invisible to every prior test, because a test built from
0-or-1-saturated color values cannot distinguish ``multiplied by one'' from ``multiplied by
two and then saturated.'' \texttt{Texture2}'s null-texture fallback had its own,
BGFX-specific bug (again, no else-branch), fixed alongside the project-wide fallback fix
above. Fog needed its own separate EasyGL fix here, because this effect uses its own
dedicated shader with no fog infrastructure at all, unlike \cnaclass{AlphaTestEffect}'s
shared shader path; Vulkan and BGFX were already covered by the project-wide fog fix.
\texttt{VertexColorEnabled} shares \cnaclass{AlphaTestEffect}'s open gap exactly: a total
no-op on both Vulkan and BGFX, for the identical reason (no color vertex attribute declared
in either backend's dual-texture shader).

\subsection{A worked example: a lightmapped floor}

\cnaclass{DualTextureEffect}'s two independent texture slots are the classic pre-baked-lighting
pattern from the original Xbox 360/Windows XNA era: a base color texture plus a second,
low-resolution lightmap multiplied over it, with no runtime lighting computation needed at all:

\begin{lstlisting}
DualTextureEffect effect(getGraphicsDeviceProperty());
effect.setWorldProperty(floorWorld);
effect.setViewProperty(camera.View());
effect.setProjectionProperty(camera.Projection());

effect.setTextureProperty(&floorDiffuse);   // base color, tiled across the floor mesh's own UVs
effect.setTexture2Property(&floorLightmap); // low-res baked lighting, its own second UV channel

for (EffectPass& pass : effect.getCurrentTechniqueProperty()->getPassesProperty()) {
    pass.Apply();
    getGraphicsDeviceProperty().DrawIndexedPrimitives(/* ... */);
}
\end{lstlisting}

The audit's headline finding directly changes what result this produces: real FNA's shader
doubles \texttt{Texture}'s RGB value before multiplying it against \texttt{Texture2} ---
\texttt{color.rgb *= 2} --- specifically so a lightmap value of 0.5 (mid-gray, the ``neutral,
unlit'' value most lightmap bakers emit) reproduces the base texture's own color exactly, rather
than darkening it to half brightness. A lightmap authored against real XNA/FNA's doubling
behavior will render visibly too dark on any CNA build predating the fix that added this factor
to all three backend shaders at once. \texttt{VertexColorEnabled} carries the identical
Vulkan/BGFX no-op gap \cnaclass{AlphaTestEffect} has, for the same underlying reason (no color
vertex attribute in either backend's dual-texture shader) --- a lightmapped floor mesh that also
tries to carry per-vertex tint runs into the same wall on those two backends.

\section{EnvironmentMapEffect}

\cnaclass{EnvironmentMapEffect} implements \texttt{IEffectMatrices},
\texttt{IEffectLights}, and \texttt{IEffectFog} together: a diffuse \texttt{Texture} plus an
\texttt{EnvironmentMap} (\cnaclass{TextureCube}, reached via the bypass mechanism described in
Chapter~\ref{ch:textures-rendertargets}), \texttt{EnvironmentMapAmount},
\texttt{EnvironmentMapSpecular}, \texttt{FresnelFactor}, \texttt{EmissiveColor}, and full
lighting. One API-fidelity detail worth noting: \texttt{LightingEnabled}'s setter
\emph{throws} if set to \texttt{false} --- matching FNA's own always-true-getter,
throws-on-false-setter behavior for this specific effect, rather than silently accepting a
value it can never actually honor. Its 14-property audit found zero property bugs, but did
surface a \texttt{Clone()} bug shared across four of the five stock effects --- cloning never
preserved \texttt{FogColor} on \cnaclass{AlphaTestEffect}, \cnaclass{DualTextureEffect},
\cnaclass{EnvironmentMapEffect}, and \cnaclass{SkinnedEffect} alike, each needing the same
one-line fix. The rendering-formula findings here are the richest of any stock effect: the
cube-map blend was additive where FNA's real formula is a \texttt{lerp}; the specular term
was added as a flat, unscaled constant rather than being scaled by the cube map's own alpha
channel (and the base blend's own color term had the identical missing-alpha-scale bug); and
--- the most striking finding --- \textbf{no Fresnel edge-weighting existed at all}, on any
backend, before this audit. Real XNA's Fresnel term (\texttt{pow(max(1 - |dot(eye, normal)|,
0), FresnelFactor) * Amount}) makes environment-map reflectivity depend on viewing angle;
CNA was applying a flat amount regardless of angle. The fix computes the term per pixel
rather than FNA's per-vertex --- a deliberate, strictly more accurate CNA deviation, not a
compatibility risk, since a per-pixel Fresnel term can only be closer to the true continuous
function than a per-vertex-interpolated approximation. A separate, genuinely interesting bug
affected the world-space normal transform on two of three backends: correct behavior under
non-uniform scale requires transforming normals by the transpose of the inverse of the
world matrix's upper $3\times3$, not the raw matrix itself; EasyGL used the raw $3\times3$,
and BGFX transformed normals with a plain matrix multiply --- both wrong the same way; Vulkan
was already correct. Both were fixed via the standard cofactor/determinant shortcut for the
transpose-inverse. All findings compose correctly in a capstone test producing byte-identical
output across all three backends.

\subsection{A worked example: a chrome car body}

A reflective, Fresnel-edge-brightened surface --- the archetypal use for this effect --- needs
both the diffuse texture and the \cnaclass{TextureCube} environment map bound, plus lighting
enabled so the Fresnel and specular terms below have a light direction to work from:

\begin{lstlisting}
EnvironmentMapEffect effect(getGraphicsDeviceProperty());
effect.setWorldProperty(carBodyWorld);
effect.setViewProperty(camera.View());
effect.setProjectionProperty(camera.Projection());

effect.setTextureProperty(&carBodyPaint);
effect.setEnvironmentMapProperty(&skybox);        // Chapter 11's TextureCube worked example
effect.setEnvironmentMapAmountProperty(0.85f);    // mostly reflective, not a full mirror
effect.setFresnelFactorProperty(1.0f);
effect.setLightingEnabledProperty(true);          // required -- setter throws if set false

for (EffectPass& pass : effect.getCurrentTechniqueProperty()->getPassesProperty()) {
    pass.Apply();
    getGraphicsDeviceProperty().DrawIndexedPrimitives(/* ... */);
}
\end{lstlisting}

Every one of this effect's richest audit findings changes what the same code above actually
renders, not merely how it renders. Before the audit's Fresnel fix, \texttt{FresnelFactor} above
had no visible effect at all --- reflectivity was flat and angle-independent regardless of its
value, rather than brightening at grazing angles the way a real Fresnel term (and real chrome)
does. Before the cube-map-blend fix, \texttt{EnvironmentMapAmount} composited additively rather
than via FNA's real \texttt{lerp}, meaning values approaching 1.0 could blow out to
over-bright rather than smoothly approaching a pure reflection. And on a non-uniformly-scaled
car body mesh specifically (common for a squashed/stretched proportions pass), the pre-fix
world-space normal transform on EasyGL and BGFX would have visibly warped the reflection's
apparent shape, because both backends transformed normals with the raw world matrix instead of
its transpose-inverse.

\section{SkinnedEffect}

\cnaclass{SkinnedEffect} caps its bone palette at \texttt{MaxBones = 72} (a static constant
that needed an out-of-line definition fix at one point for correct linkage), and exposes
\texttt{SetBoneTransforms} / \texttt{GetBoneTransforms(count)}, a \texttt{WeightsPerVertex}
property restricted to 1, 2, or 4 (throwing otherwise), \texttt{SpecularColor}/%
\texttt{SpecularPower}, full lighting and fog, and one \texttt{NOXNA} addition real XNA's
\cnaclass{SkinnedEffect} does not have at all: \texttt{VertexColorEnabled}, added
specifically to support glTF-imported meshes carrying a \texttt{COLOR\_0} vertex attribute.
The core skinning math itself --- identity, single-bone, and two-bone weighted blending ---
was pixel-verified correct on all three backends with zero bugs found; the same
\texttt{Clone()}-drops-a-field bug pattern from \cnaclass{EnvironmentMapEffect} recurred here
too, this time silently resetting \texttt{SpecularColor} / \texttt{SpecularPower} to zero on
every clone. Three further gaps were opened and closed together: \texttt{DirectionalLight1}/%
\texttt{DirectionalLight2} were unforwarded; \texttt{SpecularColor} / \texttt{SpecularPower}
had no GPU implementation on any backend at all (fixed with a real Blinn-Phong half-vector
specular term); and \texttt{WeightsPerVertex} was a complete GPU no-op --- the skinning
shader always summed all four bone weights regardless of what the property said, now fixed
with real GPU-side gating for the 2-and-4-weight cases.

\subsection{A worked example: driving SkinnedEffect from AnimationPlayer}

Chapter~\ref{ch:models-meshes}'s own worked \cnaclass{AnimationPlayer} example ended with
exactly the call this effect exists to receive --- \texttt{SetBoneTransforms} takes
\texttt{GetSkinTransforms()}'s output directly, already pre-multiplied by each bone's inverse
bind pose, which is specifically the form the GPU skinning shader expects:

\begin{lstlisting}
SkinnedEffect effect(getGraphicsDeviceProperty());
effect.setViewProperty(camera.View());
effect.setProjectionProperty(camera.Projection());
effect.setWeightsPerVertexProperty(2);    // must be 1, 2, or 4 -- throws otherwise
effect.setSpecularColorProperty(Vector3(0.3f, 0.3f, 0.3f));
effect.setSpecularPowerProperty(16.0f);
effect.EnableDefaultLighting();

// Every frame, after player.Update(...):
effect.SetBoneTransforms(player.GetSkinTransforms());
\end{lstlisting}

\texttt{MaxBones = 72} is a hard cap worth checking against your own rig before relying on this
path: a character rig with more than 72 real skinning joints (uncommon for a game character, not
unheard of for a complex creature or a converted film-asset rig) cannot be driven through this
effect at all without re-partitioning the mesh into multiple draws, each addressing a disjoint
72-bone subset. \texttt{WeightsPerVertex} is worth setting deliberately rather than leaving at
its default: before the audit's fix, this property was a complete GPU no-op --- the skinning
shader always summed all four bone weights regardless of what the property said, so a
two-weight-per-vertex rig authored assuming the third and fourth weights were ignored would
have rendered using all four anyway. Fixed, the property now genuinely gates the GPU-side sum,
which also means a rig relying on the old all-four-weights behavior while requesting
\texttt{WeightsPerVertex(2)} would see a real, visible skinning change after upgrading past the
fix --- worth flagging explicitly to anyone porting a project that predates it.

\section{The pattern behind these bugs}

Read together, these five audits tell one consistent story rather than five unrelated ones.
Almost every serious bug found falls into one of three shapes: a formula genuinely missing a
term FNA's real shader includes (the doubled-RGB dual-texture blend, the missing Fresnel
term, the missing specular implementation); a flag accepted by the public API but never
wired to any actual GPU state on one or more backends (\texttt{VertexColorEnabled} on
Vulkan/BGFX in two different effects, \texttt{WeightsPerVertex}); or a bug shared identically
across all three hardware backends because the shared, backend-agnostic C\texttt{++} code
that dispatches to per-backend shaders had the bug in the dispatch layer itself, not in any
one backend's shader (the \texttt{DirectionalLight0.Enabled} check, the project-wide fog
gap, the \texttt{Clone()}-drops-\texttt{FogColor} bug). None of these five effects' core
per-vertex/per-pixel math was ever found to be simply wrong from first principles --- every
real bug was a specific, nameable omission, which is exactly the kind of claim
Chapter~\ref{ch:what-is-cna} promised this book would keep specific rather than round off.
